%\section{事件详细过程}
\setlength{\TPHorizModule}{30mm}
\setlength{\TPVertModule}{\TPHorizModule}
%\setlength{\TPVertModule}{40mm}
\begin{changemargin}{-1cm}{-1cm}
\begin{sidewaystable} % <-- HERE
%\begin{tikzpicture}[thick,snake=zigzag, line before snake = 5mm, line after snake = 5mm]
\begin{tikzpicture}[thick, decoration={zigzag, pre length=5mm, post length=5mm}]
%draw arrow
\draw [green, ->          ] (0,8.0) -- (1,8.0);
\draw [blue,  -stealth    ] (0,7.5) -- (1,7.5) node [right] {stealth};
\draw [red,   -latex      ] (0,7.0) -- (1,7.0) node [right] {latex};
\draw [cyan,  -to         ] (0,6.5) -- (1,6.5) node [right] {to};
\draw [brown, -triangle 60] (0,6.0) -- (1,6.0) node [right] {triangle 60};
%draw dashed lines
\draw [dashed] (10,1) -- (10,11);

     % draw horizontal line   
    %  \draw (-1,0) -- (22,0); 
       \draw (-0.5,0) -- (2,0);     
      % \draw[snake] (2,0) -- (4,0);
        \draw [decorate, decoration={snake}]  (2,0) -- (4,0);
       \draw (4,0) -- (12,0);          
       \draw [decorate, decoration={snake}]  (12,0) -- (14,0);  
       \draw (14,0) -- (25.5,0);                 
      % draw vertical lines
       \foreach \x in {0,1,2,4,6,17,22,24,25.5}
       \draw (\x cm,3pt) -- (\x cm,-3pt);
        % draw nodes
       \draw (0,0) node[below=3pt] {白马} node[above=3pt] {};
          \draw (0,0) node[below=3pt] {} node[above=3pt] {印};
          \draw (0,0) node[below=3pt] {} node[above=13pt] {一};
          \draw (-0.0,0) node[below=16pt] {婦} node[above=10pt] {};
          \draw (-0.0,0) node[below=27pt] {人} node[above=10pt] {};
          \draw (-0.0,0) node[below=38pt] {懷} node[above=10pt] {};
          \draw (-0.0,0) node[below=49pt] {孕} node[above=10pt] {};
          
         \draw (-0.0,-2) node[below=60pt] {婦} node[above=10pt] {};
          \draw (-0.0,-2) node[below=71pt] {人} node[above=10pt] {};
          \draw (-0.0,-2) node[below=82pt] {懷} node[above=10pt] {};
          \draw (-0.0,-2) node[below=93pt] {孕} node[above=10pt] {};
     


         
        \draw (1,0) node[below=3pt] {} node[above=3pt] {印};
           \draw (1,0) node[below=3pt] {} node[above=13pt] {七};
           
        \draw (2,0) node[below=3pt] {} node[above=3pt] {号};
           \draw (2,0) node[below=3pt] {} node[above=13pt] {一};
           \draw (2,0) node[below=14pt] {龍} node[above=3pt] {};
           \draw (2,0) node[below=25pt] {开} node[above=3pt] {};
           \draw (2,0) node[below=36pt] {始} node[above=3pt] {};
           \draw (2,0) node[below=47pt] {摔} node[above=3pt] {};
         
         \draw (3,0) node[below=14pt] {尾巴拖拉} node[above=3pt] {};
         \draw (3,0) node[below=25pt] {三分之一} node[above=3pt] {};
         \draw (3,0) node[below=36pt] {摔到地上} node[above=3pt] {};
                 
        \draw (4,0) node[below=3pt] {} node[above=3pt] {号};
         \draw (4,0) node[below=3pt] {} node[above=13pt] {五};     
            \draw (4,0) node[below=14pt] {龍} node[above=3pt] {};
           \draw (4,0) node[below=25pt] {摔} node[above=3pt] {};
           \draw (4,0) node[below=36pt] {到} node[above=3pt] {};
           \draw (4,0) node[below=47pt] {地} node[above=3pt] {};
           \draw (4,0) node[below=58pt] {上} node[above=3pt] {};     
           
         \draw (5,0) node[below=14pt] {龍开始} node[above=3pt] {};
         \draw (5,0) node[below=24pt] {逼迫教会} node[above=3pt] {}; 
                 
       \draw (6,0) node[below=3pt] {} node[above=3pt] {号};
        \draw (6,0) node[below=14pt] {教} node[above=13pt] {七};     
        \draw (6,0) node[below=25pt] {会} node[above=13pt] {};  
        \draw (6,0) node[below=36pt] {被} node[above=13pt] {};      
        \draw (6,0) node[below=47pt] {提} node[above=13pt] {};  
 
       \draw (7.5,0) node[below=14pt] {以色列被} node[above=3pt] {};
       \draw (7.5,0) node[below=25pt] {保护三年半} node[above=3pt] {};                                
       \draw (7.5,0) node[below=40pt] {龍逼迫妇人} node[above=3pt] {};
       \draw (7.5,0) node[below=51pt] {其他的儿女} node[above=3pt] {};
       
       \draw (8,0) node[below=61pt] {} node[above=3pt] {獸};                   
       \draw (8,0) node[below=72pt] {} node[above=13pt] {中};                                       
       \draw (8,0) node[below=72pt] {} node[above=23pt] {海};                                              
        \draw (8.5,0) node[below=61pt] {} node[above=3pt] {獸};                   
       \draw (8.5,0) node[below=72pt] {} node[above=13pt] {中};                                       
       \draw (8.5,0) node[below=72pt] {} node[above=23pt] {地};                                              

       \draw (9,0) node[below=61pt] {} node[above=3pt] {羊};                   
       \draw (9,0) node[below=72pt] {} node[above=13pt] {羔};                                       
       \draw (9,0) node[below=72pt] {} node[above=23pt] {山};                                               
        \draw (9,0) node[below=61pt] {} node[above=33pt] {安};                   
       \draw (9,0) node[below=72pt] {} node[above=43pt] {錫};                                       
    
    
        \draw (11,0) node[below=61pt] {} node[above=3pt] {音};                   
       \draw (11,0) node[below=72pt] {} node[above=13pt] {声};                                       
       \draw (11,0) node[below=72pt] {} node[above=23pt] {上};                                               
        \draw (11,0) node[below=61pt] {} node[above=33pt] {天};      
   
       \draw (13,0) node[below=61pt] {} node[above=3pt] {割};                   
       \draw (13,0) node[below=72pt] {} node[above=13pt] {收};                                       
       \draw (13,0) node[below=72pt] {} node[above=23pt] {子};                                               
        \draw (13,0) node[below=61pt] {} node[above=33pt] {人};      
   
 
       \draw (17,0) node[below=3pt] {天} node[above=3pt] {碗};                   
       \draw (17,0) node[below=14pt] {使} node[above=13pt] {一};                                       
      \draw (17,0) node[below=25pt] {收} node[above=23pt] {};                                               
        \draw (17,0) node[below=36pt] {割} node[above=33pt] {};   
        
              
       \draw (20,0) node[below=3pt] {哈} node[above=3pt] {碗};                                               
       \draw (20,0) node[below=14pt] {米} node[above=13pt] {六};        
       \draw (20,0) node[below=25pt] {吉} node[above=3pt] {};                                               
       \draw (20,0) node[below=36pt] {多} node[above=13pt] {};     
       \draw (20,0) node[below=47pt] {顿} node[above=3pt] {};                                               
     
       \draw (22,0) node[below=3pt] {大} node[above=3pt] {碗};                                               
       \draw (22,0) node[below=14pt] {城} node[above=13pt] {七};     
        \draw (22,0) node[below=25pt] {地} node[above=3pt] {};                                               
       \draw (22,0) node[below=36pt] {震} node[above=13pt] {};   

      \draw (22.5,0) node[below=3pt] {巴} node[above=3pt] {};                                               
       \draw (22.5,0) node[below=14pt] {比} node[above=13pt] {};        
       \draw (22.5,0) node[below=25pt] {伦} node[above=3pt] {};                                               
       \draw (22.5,0) node[below=36pt] {倾} node[above=13pt] {};     
        \draw (22.5,0) node[below=47pt] {倒} node[above=13pt] {};    
 
 
        \draw (23,0) node[below=3pt] {羔} node[above=3pt] {};                                               
       \draw (23,0) node[below=14pt] {羊} node[above=13pt] {};        
       \draw (23.,0) node[below=25pt] {婚} node[above=3pt] {};                                               
       \draw (23.,0) node[below=36pt] {宴} node[above=13pt] {};     

        \draw (23.5,0) node[below=3pt] {白} node[above=3pt] {};                                               
       \draw (23.5,0) node[below=14pt] {马} node[above=13pt] {};        
       \draw (23.5,0) node[below=25pt] {得} node[above=3pt] {};                                               
       \draw (23.5,0) node[below=36pt] {胜} node[above=13pt] {};     
 
             
       \draw (24,0) node[below=61pt] {} node[above=3pt] {年};                   
       \draw (24,0) node[below=72pt] {} node[above=13pt] {禧};                                       
       \draw (24,0) node[below=72pt] {} node[above=23pt] {千};                                               

 
      \draw (25.5,0) node[below=3pt] {} node[above=3pt] {啟21};             
       \draw (25.5,0) node[below=3pt] {} node[above=13pt] {地};
       \draw (25.5,0) node[below=3pt] {} node[above=23pt] {新};
       \draw (25.5,0) node[below=3pt] {} node[above=33pt] {天};       
        \draw (25.5,0) node[below=3pt] {} node[above=43pt] {新};             
\end{tikzpicture}
\end{sidewaystable} 
\end{changemargin}

\newpage





\iffalse
\begin{timeline}{1}{80}{2cm}{3cm}{12cm}{18cm}{62cm}{80cm}
%\plainentry{0}{The October Revolution}
\entry{1}{七印时期}
\entry{1}{第一印：騎白馬的勝了又要勝，婦人懷孕（12：2）}
%\plainentry{2}{The October Revolution}
\plainentry{2}{七号时期}
\entry{2}{第一到五号：大紅龍尾巴拖拉著天上星辰的三分之一，开始往地上摔}
\entry{3}{第五号：“我就看見一個星從天落到地上（9：1）”。大紅龍尾巴拖拉著天上星辰的三分之一，摔在地上（12：3）}
\entry{5}{第五到七号之间：龍要吞吃婦人将生的孩子。（12：4）地上教会开始受逼迫。天上就有了爭戰，龍和牠的使者一同被摔下去。}
\entry{7}{第七号：婦人生了一個男孩子，是將來要用鐵杖轄管萬國的；她的孩子被提到神寶座那裡去了（12：5）--教会被提。}
\plainentry{8}{七号与七碗之间}
%\entry{5}{七号与七碗之间}
\entry{12}{龍逼迫生男孩子的婦人。 婦人在曠野被養活一載，二載，半載。龍與婦人其餘的兒女爭戰（12：13）。海中獸任意而行42個月。與聖徒爭戰得勝；制伏各族、各民、各方、各國。地中獸如羊羔，迷惑地上的人，叫不拜獸像的人被殺害。}
%\entry{9}{海中獸任意而行四十二個月。與聖徒爭戰得勝；制伏各族、各民、各方、各國}
%\entry{9}{地中獸如同羊羔，迷惑住在地上的人，叫所有不拜獸像的人都被殺害。}
\entry{15}{羔羊站在錫安山，十四萬四千人在寶座前，並在四活物和眾長老前唱歌}

\entry{17}{一位天使傳福音給住在地上各國、各族、各方、各民；第二位天使宣告巴比倫大城傾倒；第三位天使鼓励聖徒忍耐，守神誡命和耶穌真道}
%\entry{11}{第二位天使接著說：「叫萬民喝邪淫、大怒之酒的巴比倫大城傾倒了！傾倒了！」}
%\entry{11}{第三位天使鼓励聖徒忍耐，守神誡命和耶穌真道}

%\plainentry{7}{The October Revolution}

\entry{19}{雲上坐著的人子，手裡拿著快鐮刀。 收割地上已經熟透了的莊稼。又有一位天使從天上的殿中出來，也拿著快鐮刀收取地上葡萄樹的果子}
%\entry{12}{又有一位天使從天上的殿中出來，也拿著快鐮刀收取地上葡萄樹的果子}
%\plainentry{14}{Stock Market Crash of 1929}
%\plainentry{15}{Stock Market Crash of 1929}

\entry{41}{七碗}
%\plainentry{21}{D-Day}


\entry{44}{第六位天使把碗倒在幼發拉底大河上，我又看見三個汙穢的靈，從龍口、獸口並假先知的口中出來。 出去到普天下眾王那裡，叫他們在神全能者的大日聚集爭戰。 那三個鬼魔便叫眾王聚集在一處，希伯來話叫做哈米吉多頓。}


\entry{48}{玻璃海中勝了獸和獸的像並牠名字數目的人，唱神僕人摩西的歌和羔羊的歌}

\entry{50}{群眾在天上大聲說：哈利路亞 }
\entry{54}{羔羊婚筵}
\entry{58}{騎白馬上的显现，天上的眾軍騎著白馬跟隨他。有利劍從他口中出來，可以擊殺列國。他必用鐵杖轄管他們，並要踹全能神烈怒的酒榨。}
\plainentry{75}{撒旦被捆綁一千年，未受獸印記的同基督做王一千年}

\entry{80}{千禧年}
%\entry{23}{The Moon Landing}
\entry{86}{那一千年完了，撒旦必從監牢裡被釋放}
%\entry{98}{女人和七頭十角獸與羔羊爭戰，羔羊必勝過他們。十角與獸必恨這淫婦，使她冷落赤身，又要吃她的肉，用火將她燒盡。在一天之內，她的災殃要一齊來到}
\plainentry{90}{永世}
\entry{110}{白色的大寶座的审判}
\entry{120}{新天新地}
\entry{121}{女人和七頭十角獸與羔羊爭戰，羔羊必勝過他們。十角與獸必恨這淫婦，使她冷落赤身，又要吃她的肉，用火將她燒盡。在一天之內，她的災殃要一齊來到，就是死亡、悲哀、饑荒。她又要被火燒盡了}
\end{timeline}
\fi



